% from ega4-11-fr.tex
\oldpage[IV-3]{125}

tel que $M_{\mathfrak p}$ soit un $A$-module plat~; il existe alors
$\xi\in H$ tel que, si $\mathfrak p_\xi$ est l'image réciproque de
$\mathfrak p$ dans $B_\xi$, $(M_\xi)_{\mathfrak p_\xi}$ soit un
$A_\xi$-module plat. Soit $\xi=(\lambda,C_\lambda)$, de sorte que
l'injection $A_\xi=C_\lambda\to A_\lambda$ donne un homomorphisme
$B_\xi=B_0\otimes_{A_0}C_\lambda\to B_\lambda=B_0\otimes_{A_0}A_\lambda$,
et si $\mathfrak p_\lambda$ est l'image réciproque de $\mathfrak p$ dans
$B_\lambda$, $\mathfrak p_\xi$ est l'image réciproque de
$\mathfrak p_\lambda$ dans $B_\xi$~; on a par suite
\[
(M_\lambda)_{\mathfrak p_\lambda}
=(M_\xi\otimes_{C_\lambda}A_\lambda)_{\mathfrak p_\lambda}
=((M_\xi)_{\mathfrak p_\xi}\otimes_{C_\lambda}A_\lambda)_{\mathfrak p_\lambda},
\]
donc $(M_\lambda)_{\mathfrak p_\lambda}$ est un $A_\lambda$-module plat
($\mathbf 0_{\mathrm I}$, 6.2.1 et 6.3.2). On traite de même le cas (ii) de
l'énoncé. Nous pouvons donc par la suite supposer les $A_\lambda$
noethériens pour $\lambda\in L$ (mais non nécessairement $A$ lui-même).

\medskip
\noindent\textbf{II)} \emph{Réduction de l'énoncé global (ii) à l'énoncé
ponctuel (i).} Supposons que $\mathcal F$ soit $f$-plat. Pour tout $\lambda$,
soit $U_\lambda$ l'ensemble des $x_\lambda\in X_\lambda$ tels que
$\mathcal F_\lambda$ soit $f_\lambda$-plat au point $x_\lambda$~; on sait
que $U_\lambda$ est \emph{ouvert} dans $X_\lambda$ puisque $S_\lambda$ est
noethérien et $f_\lambda$ de type fini (11.1.1)~; soit
$V_\lambda=v_\lambda^{-1}(U_\lambda)$ son image réciproque dans $X$. Comme,
par hypothèse, pour tout $x\in X$, il y a un $\lambda$ tel que
$\mathcal F_\lambda$ soit $f_\lambda$-plat au point $x_\lambda$, projection
de $x$ dans $X_\lambda$, cela signifie que $x\in V_\lambda$ pour un
$\lambda$~; autrement dit, $X$ est réunion des $V_\lambda$. D'ailleurs
(2.1.4), pour $\lambda\leq\mu$, on a $V_\lambda\subset V_\mu$, donc, comme
$X$ est quasi-compact, il existe un indice $\mu$ tel que $X=V_\mu$. Comme
les $X_\lambda$ sont quasi-compacts, il résulte de (8.3.4) qu'il existe un
indice $\nu\geq\mu$ tel que $v_{\mu\nu}^{-1}(U_\mu)=X_\nu$~; mais par
(2.1.4), cela entraîne que $\mathcal F_\nu$ est $f_\nu$-plat.

\medskip
\noindent\textbf{III)} \emph{Fin de la démonstration.} Il reste à prouver
(i) en supposant $S_0$ affine et noethérien~; si $y_0$ est la projection de
$x$ dans $S_0$, on peut en outre, en vertu de (2.1.4) et de
($\mathbf I$, 3.6.5), remplacer $S_0$ par
$\operatorname{Spec}(\mathcal O_{y_0})$, autrement dit on peut se borner au
cas où $A_0$ est un anneau \emph{local} noethérien, d'idéal maximal
$\mathfrak m_0=\mathfrak j_{y_0}$~; par définition, $\mathfrak m_0$ est
l'image réciproque de l'idéal premier $\mathfrak p=\mathfrak j_x$ de $B$,
et $M_{\mathfrak p}$ est supposé être un $A$-module plat~; on a donc en
particulier $\operatorname{Tor}_1^A(A/\mathfrak m_0A,M_{\mathfrak p})=0$,
ce qui s'écrit aussi, puisque les $\operatorname{Tor}_1^A(A/\mathfrak m_0A,M)$
sont des $B$-modules et que $B_{\mathfrak p}$ est plat sur $B$,
$(\operatorname{Tor}_1^A(A/\mathfrak m_0A,M))_{\mathfrak p}=0$.
Notons maintenant que $\operatorname{Tor}_1^{A_0}(A_0/\mathfrak m_0,M_0)$
est un $B_0$-module \emph{de type fini}, car on peut le définir en prenant
une résolution de $A_0/\mathfrak m_0$ par des $A_0$-modules libres
\emph{de type fini} ($A_0$ étant noethérien) et en tensorisant par $M_0$, ce
qui donne des $B_0$-modules de type fini~; comme $B_0$ est noethérien,
l'homologie du complexe ainsi obtenu est bien formée de $B_0$-modules de
type fini. Soit $(t_i^0)_{1\leq i\leq m}$ un système de générateurs du
$B_0$-module $\operatorname{Tor}_1^{A_0}(A_0/\mathfrak m_0,M_0)$ et soit
$t_i$ l'image canonique (11.2.1.1) de $t_i^0$ dans
$\operatorname{Tor}_1^A(A/\mathfrak m_0A,M)$. L'hypothèse entraîne qu'il
existe un $h\in B-\mathfrak p$ tel que $ht_i=0$ pour $1\leq i\leq m$. Or,
on a (11.2.2.1)
\[
\operatorname{Tor}_1^A(A/\mathfrak m_0A,M)=
\varinjlim\operatorname{Tor}_1^{A_\lambda}
(A_\lambda/\mathfrak m_0A_\lambda,M_\lambda)~;
\]
il existe donc un $\lambda$ tel que, si les $t_i^\lambda$ sont les images
des $t_i^0$ dans
$\operatorname{Tor}_1^{A_\lambda}(A_\lambda/\mathfrak m_0A_\lambda,
M_\lambda)$, il existe $h_\lambda\in B_\lambda$ d'image $h$, tel que
$h_\lambda t_i^\lambda=0$ pour $1\leq i\leq m$. Soit
$\mathfrak p_\lambda=\mathfrak j_{x_\lambda}$ l'idéal premier de
$B_\lambda$ image réciproque de $\mathfrak p$~; on a
$h_\lambda\in B_\lambda-\mathfrak p_\lambda$, donc les images canoniques
des $t_i^0$ dans
$(\operatorname{Tor}_1^{A_\lambda}(A_\lambda/\mathfrak m_0A_\lambda,
M_\lambda))_{\mathfrak p_\lambda}$ sont nulles et par suite
l'homomorphisme
$\operatorname{Tor}_1^{A_0}(A_0/\mathfrak m_0,M_0)\to
(\operatorname{Tor}_1^{A_\lambda}(A_\lambda/\mathfrak m_0A_\lambda,
M_\lambda))_{\mathfrak p_\lambda}$ est nul. Les conditions du lemme
(11.2.5) sont donc remplies ($A_0/\mathfrak m_0$ étant un corps), et
$(M_\lambda)_{\mathfrak p_\lambda}$ est un $A_\lambda$-module plat, ce qui
achève de démontrer le théorème.

\begin{corollary}[11.2.7]
\label{IV.11.2.7-fr}
Soient $S=\operatorname{Spec}(A)$ un schéma affine, $f:X\to S$ un morphisme,
$\mathcal F$ un $\mathcal O_X$-Module quasi-cohérent, $x$ un point de $X$.
Les conditions suivantes sont équivalentes~:
\begin{enumerate}[label=\emph{\alph*)}]
\item $f$ est un morphisme de présentation finie, $\mathcal F$ est un
$\mathcal O_X$-Module de présentation finie et $\mathcal F$ est $f$-plat au
point $x$ (resp. $f$-plat).

